\section{Scope, antecedents, and principal results}
\label{intro:sec:scope}

The starting point is the manuscript in
\path{Analysis/Polylogarithms/docs/manuscript} of the ProveIt repository,
at commit \texttt{a2a4cf58c49c745058c40e4a6748d472a3420f18}
\cite{ProveItSnapshot}. The associated incoming reports were inspected
at the same commit. This matters mathematically: several questions that
still appear open in older chapter text have already been answered in
those reports. In particular, the Gaussian weight-five and weight-six
reductions, the residual level-three and level-four polynomial spaces,
and the Möbius-duality certificate for $S_4$ are antecedents, not results
newly established here \cite{PriorResearch,PriorRigidity}.

The continuation develops four connected directions. Each replaces a
finite observation or a qualitative boundary statement by an exact
structural result. Polynomial symmetry determines an infinite family of
matrix ranks. Exterior symbols determine the formal reduction type of
every positive rational Herglotz argument. A singular perturbation
analysis explains why real zeros disappear near the elementary Lerch
endpoint and why the surviving branches need not be monotone. Finite
Lagrange inversion gives the entire exponential scale of an angular
zero, while a separate argument shows that the resulting infinite
expansion is necessarily divergent.

\subsection{Four different meanings of an exact assertion}

We use the following distinctions throughout.
\begin{enumerate}
\item An \emph{analytic identity} is an equality of specified functions
or numbers, with domains, branches, constants, and limiting operations
fixed. The explicit formulas for $J(3/5)$ and $J(n/(n+1))$ have this
meaning.
\item A \emph{coefficient-rank theorem} concerns a fully specified
rational matrix generated by shuffle and stuffle rows. Its kernel is
not a space of numerical periods. A complete answer for that matrix
does not exclude identities arising from additional relations.
\item A \emph{formal dilogarithm classification} concerns the
rationalized pre-Bloch group and its boundary map. A nonzero obstruction
excludes a reduction in that calculus. It does not prove that a
numerical coincidence between periods is impossible.
\item An \emph{exact certificate} is a finite rational or integer
calculation; when it certifies an infinite analytic expression, it is
accompanied by a stated error theorem. Floating-point root
plots and high-precision identity checks are supplementary diagnostics;
they are never used as substitutes for the analytic proofs.
\end{enumerate}
These distinctions allow strong completeness statements without
silently assuming transcendence or period conjectures.

\subsection{The rank conjecture and its universal extension}

The source's Conjecture \texttt{signed:conj:rank} uses $6w-7$ imaginary
depth-two coordinates at fourth roots of unity. It removes the divergent
coordinate and retains a single regularized difference at the boundary.
Theorem~\ref{fg:thm:rank} proves both proposed odd-weight ranks and also
gives the even-weight answer. More significantly, Theorem
\ref{fg:thm:kernel} describes the entire kernel through two freely
chosen homogeneous polynomials. The previously unresolved issue is the
canonical lift of the same-point polynomial through the precise
boundary convention; a rank fit in low weights cannot supply this step.

Theorem~\ref{fg:thm:universal} shows that the lift is independent of the
cyclotomic level. It splits off one explicit antireflection sector for
each conjugate pair of nonreal roots. At level three, combining this
new lift with the earlier residual invariant ring gives the full ranks.
The results provide exact separating functionals for proposed rows,
as well as completeness of the same-point rows in the stated system.
They should be compared with the broader standard relations studied
by Zhao \cite{Zhao2010} and with parity reductions
\cite{Panzer2017,Umezawa2025}; none of those broader theories is being
identified with this restricted single-product matrix.

\subsection{All rational arguments of the Herglotz companion}

Let
\[
J(x)=\int_0^1\frac{\log(1+t^x)}{1+t}\,dt,\qquad x>0.
\]
For reduced $p/q$ with both entries odd, Theorem
\ref{jnew:thm:classification} proves that a formal rational-dilogarithm
reduction exists precisely when $p,q\in\{1,3,5\}$. A formal logarithmic
reduction in this parity class occurs only at $p=q=1$. If $e$ is the
even entry and $o$ the odd entry, the two reduction types coincide and
are characterized by
\[
e^2\equiv\pm1\pmod o,\qquad o^2\equiv1\pmod{2e}.
\]
The doubled modulus is essential. The proof adds a proper-real-subfield
faithfulness theorem for an individual cyclotomic symbol and an exact
dyadic norm identity to the source's all-conductor Gram theorem.

The analytic evaluation section fixes the logarithmic and $\pi^2$
terms suppressed by the formal notation. It derives an explicit
log-sine formula for every consecutive rational argument and obtains
\[
\begin{aligned}
J(3/5)={}&\Li_2(1/3)-\tfrac12\Li_2(1/5)-\tfrac{7\pi^2}{180}
 +\tfrac12\log^2 2+\tfrac12\log^2 3\\
&-\tfrac14\log^2 5+\log^2\left(\tfrac{1+\sqrt5}{2}\right).
\end{aligned}
\]
The functional equations of Radchenko--Zagier and Choie--Kumar are
credited explicitly \cite{RadchenkoZagier,ChoieKumar2023}. The finite
difference used here is a reformulation of known functional equations,
not a claimed new general functional equation. The classifications and
explicit specializations are the contributions of this continuation.

\subsection{Zero geometry at both Lerch endpoints}

For $n,k\geq1$, the normalized differentiated logarithmic Lerch sum is
\[
G_{n,k}(a,\rho)=\sum_{m\geq0}\rho^m
\frac{P_{n,k}(-\log(a+m))}{(a+m)^{k+1}},\qquad a>0,\quad 0\leq\rho\leq1.
\]
The polynomial $P_{n,k}$ is explicit. At $\rho=0$ it may have a
multiple zero at $a=1$. Theorem~\ref{lerchboundary:thm:smallrho}
resolves that zero into convergent Puiseux branches and gives the exact
number of positive zeros for every fixed $n,k$ and all sufficiently
small positive $\rho$. A common compact zero region proves that the
local analysis captures every positive zero.

For $n=2$ we go further: there are exactly two simple positive zeros
for every $k\geq1$ and every $\rho\in[0,1]$. Thus the source's uniform
saturation threshold has the sharp value $K_2=1$. At $k=1$ the smaller
zero first moves below $1$ and ends above $1$, so it is nonmonotone.
The larger branch is strictly decreasing; the smaller branch crosses
$1$ exactly once in the interior and increases thereafter.
At the opposite endpoint, a universal term
$\tau^k\log^{n+1}(1/\tau)$, with $\tau=-\log\rho$, makes every simple
endpoint zero exactly $C^{k-1}$ and not $C^k$. For the two $n=2,k=1$
branches the slopes tend respectively to $+\infty$ and $-\infty$.

\subsection{Every angular scale, and why the series diverges}

Let $\theta_{a,b}$ be the source's unique zero of
$\Imag\Li_{a,b}(e^{\ii\theta},1)$ in the upper semicircle. We construct
the expansion of $\delta_{a,b}=\pi/2-\theta_{a,b}$ to any prescribed
exponential accuracy, uniformly for all real $b>0$. Its scales are
products of the rational numbers $2/n$, $n\geq3$, and its coefficients
are exact polynomials in generalized harmonic numbers. Every cutoff
requires only finitely many coefficients. Ten explicit terms and a
sharp signed remainder are tabulated.

The final asymptotic theorem also clarifies what ``complete'' means.
At each fixed $a,b$, a prime-indexed subsequence of the grouped terms
does not even tend to zero. Consequently the infinite series cannot
be summed by ordinary convergence, although every fixed exponential
cutoff is rigorously asymptotic as $a\to\infty$. This is a concrete
example in which exact, arbitrarily long expansions coexist with
unavoidable divergence.

\subsection{Dependencies and reproducibility}

The matrix proofs and angular coefficient construction are algebraic
and analytic arguments valid in all weights. The Herglotz classification
uses the source's proved Gram theorem and its explicitly stated
rationalized Bloch-group facts. The zero results use the source's
global variation bound and eleven replayed rational endpoint signs;
the essential background and the variation argument are recalled here.
Only the nonvanishing of a polynomial at $-\log2$ uses classical
Hermite--Lindemann transcendence.

The source manifest, scripts, output receipts, and a claim ledger are
included with the article. New proofs received independent internal
mathematical cross-checks. No claim of formal proof-assistant
verification or external peer review is made. The remaining questions
are listed with concrete next steps in Section~\ref{research:sec:program}.
